покой (все детали попарно) & 1 & 0\\
челнок на всех позициях и во всех рядах, включая покой под строкой (с шагом пол-позиции) & 231 & 0\\
отвод узла 0, 9 и 17\,u при разных позициях, подъёме и ударе & 45 & 0\\
ход каретки по всей строке, молоток взведён и в ударе & 18 & 0\\
каждая клавиша нажата в обоих регистрах, при открытой и запертой гребёнке; рамки сдвинуты по коду & 196 & 0\\
оборот циклового вала через 10° для четырёх крайних знаков (с лентой, взводом, защёлками и блокировкой) & 144 & 0\\
мёртвый знак при каждом значении памяти B и C: в начале, на 40\,\% и у ступени; второй сектор в крайних положениях & 114 & 0\\
спуск: упор U для всех ширин при $\pm u$, память M, забой и «назад на литеру» в трёх местах строки & 109 & 0\\
чтение кода (ред. 0.10): каждая продвигающая литера на оси удара, щупы на полпути и на зубцах, ползунки и тяга W с ними & 144 & 0\\
холостой цикл (ред. 0.11): мёртвые знаки над и под строкой и пробел, оборот через 10°, запоры чтения и удара введены & 108 & 0\\
